% Workbench-ready, self-contained mathematical section.
% Required bibliography keys: Borovoi_2001, GorodnikOhShah_2009,
% Voight_2021.  Every cross-reference label is namespaced koide-count-*.

\section{Exact arithmetic of the Koide defect-one shell}
\label{koide-count-sec-main}

This section isolates an arithmetic result that does not use the existence of
the proposed complex structure on \(S^6\), any period-map assertion, or any
Yang--Mills interpretation.  All coordinates, chambers, local factors, and
normalizations used in the final constant are fixed explicitly below.

For \(r=(r_1,r_2,r_3)\), define the Koide quadratic form
\begin{equation}
 Q_K(r_1,r_2,r_3)
 :=\frac23(r_1+r_2+r_3)^2-(r_1^2+r_2^2+r_3^2).
 \label{koide-count-eq-qk}
\end{equation}
For integral coordinates put
\begin{equation}
 \Delta(x,y,z):=4(xy+xz+yz)-(x^2+y^2+z^2).
 \label{koide-count-eq-defect}
\end{equation}
Expanding \eqref{koide-count-eq-qk} gives the exact identity
\begin{equation}
 \Delta(x,y,z)=3Q_K(x,y,z).
 \label{koide-count-eq-defect-qk}
\end{equation}
Thus the equation studied here,
\(\Delta=-1\), is the nonzero level \(Q_K=-1/3\); it is not an affine
translate of the null cone \(Q_K=0\).

\subsection{Vieta reduction and exact ternary coordinates}
\label{koide-count-subsec-vieta}

Regarded as a monic quadratic equation in \(z\), the defect equation is
\begin{equation}
 z^2-4(x+y)z+x^2+y^2-4xy-1=0.
 \label{koide-count-eq-monic-z}
\end{equation}
Consequently the other root is
\begin{equation}
 z'=4(x+y)-z.
 \label{koide-count-eq-vieta-root}
\end{equation}
The analogous substitutions in \(x\) and \(y\) are integral involutions
preserving \(\Delta\).

Let $e_1,e_2,e_3$ be the standard basis and let $B_K$ be the integral
bilinear form with Gram matrix
\begin{equation}
 A_K=
 \begin{pmatrix}
 -1&2&2\\
 2&-1&2\\
 2&2&-1
 \end{pmatrix},
 \qquad
 B_K(u,v)=u^{\mathsf T}A_Kv.
 \label{koide-count-eq-vieta-gram}
\end{equation}
Thus $\Delta(u)=B_K(u,u)$.  The eigenvalue of $A_K$ on
$\mathbb R(1,1,1)$ is $3$, and its eigenvalue on the orthogonal complement is
$-3$, so this integral presentation has signature $(1,2)$.

\begin{proposition}[Vieta moves are Lorentzian reflections]
\label{koide-count-prop-vieta-reflections}
For $i=1,2,3$, define
\begin{equation}
 s_i(u)=u-
 \frac{2B_K(u,e_i)}{B_K(e_i,e_i)}e_i.
 \label{koide-count-eq-reflection-definition}
\end{equation}
Then $s_i\in\operatorname O(B_K,\mathbb Z)$,
$s_i^2=1$, and $\det s_i=-1$.  Explicitly,
\begin{align}
 s_1(x,y,z)&=(-x+4y+4z,y,z),\label{koide-count-eq-vieta-s1}\\
 s_2(x,y,z)&=(x,-y+4x+4z,z),\label{koide-count-eq-vieta-s2}\\
 s_3(x,y,z)&=(x,y,-z+4x+4y).\label{koide-count-eq-vieta-s3}
\end{align}
Consequently these are exactly the three Vieta root substitutions for
$\Delta=-1$, now identified as reflections of the unchanged Lorentzian
lattice $(\mathbb Z^3,B_K)$.
\end{proposition}

\begin{proof}
Since $B_K(e_i,e_i)=-1$, formula
\eqref{koide-count-eq-reflection-definition} is integral and direct
evaluation gives \eqref{koide-count-eq-vieta-s1}--
\eqref{koide-count-eq-vieta-s3}.  The standard reflection calculation
gives
\[
 B_K(s_i(u),s_i(v))=B_K(u,v),\qquad s_i^2=1.
\]
The map fixes the hyperplane $e_i^\perp$ pointwise and sends $e_i$ to
$-e_i$, hence has determinant $-1$.  Finally, when the equation is viewed
as a monic quadratic in one coordinate, the displayed replacement is the
sum of its two roots minus the original root, proving the last assertion.
\end{proof}

\begin{proposition}[Strict Vieta reduction]
\label{koide-count-prop-vieta-reduction}
Let \(0<x\leq y\leq z\) be integers satisfying
\(\Delta(x,y,z)=-1\).  If \(z<4(x+y)\), then the other root in
\eqref{koide-count-eq-vieta-root} satisfies \(0<z'<y\).  If
\(z\geq4(x+y)\), then \(z'\leq0\).
\end{proposition}

\begin{proof}
In the first case \(z'>0\) follows immediately from
\eqref{koide-count-eq-vieta-root}.  Let \(P(Z)\) be the polynomial on the
left of \eqref{koide-count-eq-monic-z}.  Since \(x\leq y\),
\[
 P(y)=x^2-8xy-2y^2-1
 \leq xy-8xy-2y^2-1<0.
\]
But \(P(y)=(y-z)(y-z')\), and \(y-z\leq0\); hence \(y-z'>0\), so
\(z'<y\).  The second assertion again follows directly from
\eqref{koide-count-eq-vieta-root}.
\end{proof}

Define the reduced-root set and its height count by
\begin{align}
 \mathcal R
 &:={}
 \{(x,y,z)\in\mathbb Z_{>0}^3:
 x\leq y\leq z,\ \Delta(x,y,z)=-1,\ z\geq4(x+y)\},
 \label{koide-count-eq-reduced-set}\\
 R(X)&:=\#\{(x,y,z)\in\mathcal R:z\leq X\}.
 \label{koide-count-eq-rx-definition}
\end{align}
Proposition~\ref{koide-count-prop-vieta-reduction} shows that the wall
\(z\geq4(x+y)\) is exactly the terminal side of the decreasing Vieta move
after the coordinates have been ordered.

\begin{proposition}[The positive Vieta graph is a rooted forest]
\label{koide-count-prop-vieta-forest}
Let $\mathcal S^+$ be the set of unordered triples of positive integers
satisfying $\Delta=-1$.  Join two vertices in a simple graph when one is
obtained from the other by a Vieta reflection $s_i$ and the reflected
coordinate is positive.  Every connected component of this graph is a
tree with exactly one root, and the roots are precisely the unordered
triples whose ordered representatives lie in $\mathcal R$.  Equivalently,
every positive integral solution has a unique reduced representative, and
it is obtained by repeatedly ordering the coordinates and reflecting the
largest coordinate whenever it lies below the terminal wall.
\end{proposition}

\begin{proof}
First, no Vieta reflection fixes a point of $\Delta=-1$.  Indeed, a fixed
reflection in a coordinate $w$, with the other two coordinates denoted by
$u,v$, would give $w=2(u+v)$, and substitution in the unchanged form gives
\[
 \Delta(u,v,2(u+v))=3(u^2+4uv+v^2)>0.
\]
Thus the sum of the three coordinates changes strictly across every edge.

Now order a vertex as $0<x\leq y\leq z$.  In fact $y<z$, because $y=z$
would give
\[
 \Delta(x,y,y)=-x^2+8xy+2y^2>0.
\]
Reflecting either of the two smaller coordinates produces a new coordinate
strictly larger than $z$:
\[
 4(y+z)-x>z,\qquad 4(x+z)-y>z.
\]
If $z<4(x+y)$, Proposition~\ref{koide-count-prop-vieta-reduction} says
that reflection in the largest coordinate produces $0<z'<y$ and strictly
decreases the coordinate sum.  If $z\geq4(x+y)$, that reflection is
nonpositive and hence is not an edge of $\mathcal S^+$.  Therefore every
nonterminal vertex has exactly one neighbour of smaller coordinate sum,
whereas a terminal vertex has none.

Iteration of the unique decreasing move terminates because the coordinate
sum is a positive integer.  The endpoint is represented by an element of
$\mathcal R$.  The graph has no cycle: on a putative simple cycle choose a
vertex of largest coordinate sum; its two neighbours on the cycle would be
two distinct smaller neighbours, contrary to uniqueness.  Thus each
component is a tree.  A tree in which following the unique decreasing edge
always terminates has exactly one terminal vertex, proving both uniqueness
and the characterization of the roots.
\end{proof}

For an ordered triple define
\begin{equation}
 n=z-4(x+y),\qquad q=y-x,\qquad a=x+y-4n,
 \label{koide-count-eq-forward-coordinates}
\end{equation}
and let
\begin{equation}
 \mathcal T:=
 \left\{(a,q,n)\in\mathbb Z^3:
 \begin{array}{c}
 a>0,\ q\geq0,\ n\geq0,\ a+4n>q,\\
 a^2-3q^2-18n^2=-2
 \end{array}
 \right\}.
 \label{koide-count-eq-ternary-chamber}
\end{equation}

\begin{theorem}[Integral ternary-coordinate bijection]
\label{koide-count-thm-coordinate-bijection}
The map in \eqref{koide-count-eq-forward-coordinates} restricts to a
bijection \(\Phi:\mathcal R\to\mathcal T\).  Its inverse is
\begin{equation}
 (a,q,n)\longmapsto
 \left(
 \frac{a+4n-q}{2},
 \frac{a+4n+q}{2},
 4a+17n
 \right).
 \label{koide-count-eq-inverse-coordinates}
\end{equation}
For arbitrary integral \((x,y,z)\) and the associated \((a,q,n)\),
\begin{equation}
 2\Delta(x,y,z)=a^2-3q^2-18n^2,
 \qquad z=4a+17n.
 \label{koide-count-eq-polynomial-identity}
\end{equation}
In particular,
\begin{equation}
 R(X)=\#\{(a,q,n)\in\mathcal T:4a+17n\leq X\}.
 \label{koide-count-eq-exact-count}
\end{equation}
\end{theorem}

\begin{proof}
Put \(p=x+y\).  Since \(q=y-x\), one has
\(2x=p-q\) and \(2y=p+q\), and direct substitution into
\eqref{koide-count-eq-defect} gives
\[
 2\Delta=p^2-3q^2+8pz-2z^2.
\]
The definition \(n=z-4p\) gives \(z=4p+n\), whence
\[
 2\Delta=p^2-3q^2-8pn-2n^2.
\]
Finally \(a=p-4n\), equivalently \(p=a+4n\), gives both identities in
\eqref{koide-count-eq-polynomial-identity}.

Suppose first that \((x,y,z)\in\mathcal R\).  Then \(q,n\geq0\), while
\(a+4n=p>q\) because \(p-q=2x>0\).  The defect equation and
\eqref{koide-count-eq-polynomial-identity} give
\(a^2-3q^2-18n^2=-2\).  If \(n=0\), then \(a=p>0\).  If \(n\geq1\),
then
\[
 a^2=3q^2+18n^2-2\geq18n^2-2\geq16n^2.
\]
Were \(a\leq0\), this inequality would force \(a\leq-4n\), contrary to
\(a+4n>q\geq0\).  Hence \(a>0\), and the forward map lands in
\(\mathcal T\).

Conversely, reducing the ternary equation modulo \(2\) gives
\(a\equiv q\pmod2\), so all three coordinates in
\eqref{koide-count-eq-inverse-coordinates} are integral.  Put
\(p=a+4n\).  The chamber inequality gives
\[
 x=\frac{p-q}{2}>0,\qquad
 y=\frac{p+q}{2}\geq x.
\]
Moreover
\[
 z=4a+17n=4p+n\geq4p>p>y.
\]
The identities \eqref{koide-count-eq-polynomial-identity} now give
\(\Delta=-1\) and \(z-4(x+y)=n\geq0\).  Thus the inverse map lands in
\(\mathcal R\).  The displayed formulae show directly that the two
compositions are the identity.  The last assertion follows from the height
identity in \eqref{koide-count-eq-polynomial-identity}.
\end{proof}

\begin{proposition}[Complete Pell classification of the chamber faces]
\label{koide-count-prop-pell-faces}
The integral points of $\mathcal T$ on its two coordinate faces are the
following two disjoint families.

On the terminal Vieta wall $n=0$, define integers $a_k,q_k$ by
\begin{equation}
 a_k+q_k\sqrt3=(1+\sqrt3)(2+\sqrt3)^k
 \qquad(k\geq0).
 \label{koide-count-eq-wall-pell-sequence}
\end{equation}
Then the points of $\mathcal T$ with $n=0$ are exactly
$(a_k,q_k,0)$ for $k\geq1$, and their reduced roots are
\begin{equation}
 \left(
  \frac{a_k-q_k}{2},
  \frac{a_k+q_k}{2},
  4a_k
 \right).
 \label{koide-count-eq-wall-pell-roots}
\end{equation}
The first is $(1,4,20)$.

On the symmetry face $q=0$, define integers $A_k,N_k$ by
\begin{equation}
 A_k+N_k\sqrt{18}=(4+\sqrt{18})(17+4\sqrt{18})^k
 \qquad(k\geq0).
 \label{koide-count-eq-symmetric-pell-sequence}
\end{equation}
Then the points of $\mathcal T$ with $q=0$ are exactly
$(A_k,0,N_k)$ for $k\geq0$, and their reduced roots are
\begin{equation}
 \left(
  \frac{A_k+4N_k}{2},
  \frac{A_k+4N_k}{2},
  4A_k+17N_k
 \right).
 \label{koide-count-eq-symmetric-pell-roots}
\end{equation}
The first is $(4,4,33)$.

In the projectivized real quadric these families tend, respectively, to
the two chamber-boundary rays
\begin{equation}
 (a,q,n)=\mathbb R_{>0}(\sqrt3,1,0),
 \qquad
 (a,q,n)=\mathbb R_{>0}(3\sqrt2,0,1).
 \label{koide-count-eq-pell-limit-rays}
\end{equation}
Each family contributes only $O(\log X)$ points to $R(X)$.
\end{proposition}

\begin{proof}
On $n=0$, the equation of $\mathcal T$ is
$a^2-3q^2=-2$.  Here $q=0$ is impossible.  If $q=1$, positivity gives
$a=1$, which fails the strict chamber inequality $a>q$.  For any positive
solution with $q\geq2$, multiplication by the norm-one unit
$2-\sqrt3$ gives
\begin{equation}
 a'+q'\sqrt3=(a+q\sqrt3)(2-\sqrt3),\qquad
 a'=2a-3q,\quad q'=2q-a.
 \label{koide-count-eq-wall-pell-descent}
\end{equation}
The unchanged equation implies
\[
 a<2q,\qquad
 a>\frac32q,\qquad
 a>q
 \quad(q\geq2),
\]
because its squares differ from $4q^2$, $9q^2/4$, and $q^2$ by,
respectively, $q^2+2$, $3q^2/4-2>0$, and $2q^2-2>0$.
Thus $a',q'>0$ and $q'<q$.  Repeated descent reaches $(1,1)$, and reversing
the descent proves \eqref{koide-count-eq-wall-pell-sequence} is the complete
positive solution set.  The recurrence
\[
 a_{k+1}=2a_k+3q_k,\qquad q_{k+1}=a_k+2q_k
\]
shows inductively that $a_k,q_k$ are both odd.  The chamber excludes only
$k=0$, and Theorem~\ref{koide-count-thm-coordinate-bijection} now gives
\eqref{koide-count-eq-wall-pell-roots}.

On $q=0$, one has $a^2-18n^2=-2$.  Necessarily $n>0$.  Direct substitution
for $1\leq n\leq5$ gives the sole solution $(a,n)=(4,1)$.  If $n\geq6$,
multiplication by the norm-one unit $17-4\sqrt{18}$ gives
\begin{equation}
 a'+n'\sqrt{18}=(a+n\sqrt{18})(17-4\sqrt{18}),\qquad
 a'=17a-72n,\quad n'=17n-4a.
 \label{koide-count-eq-symmetric-pell-descent}
\end{equation}
Both new coordinates are positive.  Indeed,
\[
 \frac an=\sqrt{18-\frac2{n^2}}
 >\frac{72}{17}\quad(n\geq6),
 \qquad
 \frac an<\sqrt{18}<\frac{17}{4}.
\]
Moreover $a>4n$ and hence $0<n'<n$.  Descent therefore reaches $(4,1)$,
and reversal proves \eqref{koide-count-eq-symmetric-pell-sequence} is the
complete positive solution set.  Its recurrence,
\[
 A_{k+1}=17A_k+72N_k,\qquad
 N_{k+1}=4A_k+17N_k,
\]
keeps $A_k$ even and $N_k$ odd, so
\eqref{koide-count-eq-symmetric-pell-roots} follows from the inverse
coordinate formula.

Conjugating the two identities defining the sequences shows
\[
 \frac{a_k}{q_k}\longrightarrow\sqrt3,\qquad
 \frac{A_k}{N_k}\longrightarrow\sqrt{18}=3\sqrt2,
\]
which proves \eqref{koide-count-eq-pell-limit-rays}.  Finally, the dominant
factors $2+\sqrt3>1$ and $17+4\sqrt{18}>1$ give exponential growth in $k$,
whereas the heights in \eqref{koide-count-eq-wall-pell-roots} and
\eqref{koide-count-eq-symmetric-pell-roots} are positive linear forms in
the displayed coordinates.  Hence either face contains only $O(\log X)$
points of height at most $X$.
\end{proof}

\subsection{The affine quadric and its local densities}
\label{koide-count-subsec-local}

Put
\begin{equation}
 F(a,q,n):=a^2-3q^2-18n^2,
 \qquad
 \mathscr X:=\{F=-2\}\subset\mathbb A^3.
 \label{koide-count-eq-affine-quadric}
\end{equation}
The matrix of \(F\) is
\(A=\operatorname{diag}(1,-3,-18)\), so
\(\det A=54\), and
\(-(-2)\det A=108\) is not a square.  The integral point
\(x_0=(1,1,0)\) lies on \(\mathscr X\).

For every prime \(p\), define
\begin{equation}
 N_{p,k}:=
 \#\{(a,q,n)\bmod p^k:F(a,q,n)\equiv-2\pmod{p^k}\},
 \qquad
 \sigma_p:=\lim_{k\to\infty}p^{-2k}N_{p,k}.
 \label{koide-count-eq-local-density-definition}
\end{equation}

\begin{theorem}[Exact local densities]
\label{koide-count-thm-local-densities}
The limits in \eqref{koide-count-eq-local-density-definition} exist and
are
\begin{equation}
 \sigma_2=\frac32,
 \qquad
 \sigma_3=2,
 \qquad
 \sigma_p=
 \frac{1-p^{-2}}{1-\left(\frac3p\right)p^{-1}}
 \quad(p\geq5).
 \label{koide-count-eq-all-local-densities}
\end{equation}
Writing \(\chi_{12}\) for the primitive real character modulo \(12\), the
singular series, with the primes taken in increasing order (equivalently,
with Borovoi's convergence-factor prescription), is therefore
\begin{equation}
 \mathfrak S(\mathscr X):=\prod_p\sigma_p
 =\frac92\frac{L(1,\chi_{12})}{\zeta(2)}
 =\frac92\frac{\log(2+\sqrt3)}{\sqrt3\,\zeta(2)}.
 \label{koide-count-eq-singular-series}
\end{equation}
\end{theorem}

\begin{proof}
We compute the two bad primes first.

For \(p=3\), once \(q,n\bmod3^k\) have been chosen, the equation for
\(a\) is
\[
 a^2\equiv-2+3q^2+18n^2\equiv1\pmod3.
\]
The two roots modulo \(3\) are nonsingular, and each has a unique Hensel
lift modulo every \(3^k\).  Thus
\begin{equation}
 N_{3,k}=2\cdot3^{2k},
 \label{koide-count-eq-three-adic-count}
\end{equation}
which gives \(\sigma_3=2\).

For \(p=2\), reduction modulo \(4\) leaves exactly two parity branches:
\begin{equation}
 (a,q,n)\equiv(0,0,1)\pmod2,
 \qquad
 (a,q,n)\equiv(1,1,0)\pmod2.
 \label{koide-count-eq-two-adic-branches}
\end{equation}
Fix \(k\geq4\) and put \(j=k-1\).  On the first branch write
\(a=2A\), \(q=2Q\), and put \(r=3n\).  Then
\begin{equation}
 F(a,q,n)+2
 =-2\bigl(r^2-2A^2+6Q^2-1\bigr).
 \label{koide-count-eq-two-adic-even-branch}
\end{equation}
Hence we must solve
\[
 r^2\equiv1+2A^2-6Q^2\pmod{2^j}.
\]
An odd unit is a square modulo \(2^j\), \(j\geq3\), precisely when it is
\(1\pmod8\), in which case it has four roots.  Here that condition is
equivalent to \(A,Q\) both being even.  There are \(2^{2j-2}\) such pairs
\((A,Q)\bmod2^j\), four choices of \(r\bmod2^j\), and, because \(3\) is
a unit, two odd choices of \(n\bmod2^k\) above each such \(r\).  This
branch therefore contributes
\begin{equation}
 2^{2j-2}\cdot4\cdot2=2^{2k-1}
 \label{koide-count-eq-two-adic-first-branch-count}
\end{equation}
solutions modulo \(2^k\).

On the second branch write \(n=2m\) and set
\begin{equation}
 v=\frac{a-q}{2},
 \qquad
 r=\frac{-a+3q}{2}.
 \label{koide-count-eq-two-adic-divided-coordinates}
\end{equation}
Then \(r,v\) have opposite parity, the inverse formula is
\(a=r+3v\), \(q=r+v\), and
\begin{equation}
 F(a,q,n)+2
 =-2\bigl(r^2-3v^2+36m^2-1\bigr).
 \label{koide-count-eq-two-adic-odd-branch}
\end{equation}
First suppose \(v=2V\) and \(r\) is odd.  For fixed
\(V\bmod2^{j-1}\) and \(m\bmod2^j\), the equation for \(r\) is
\[
 r^2\equiv1+12V^2-36m^2\pmod{2^j}.
\]
Its right side is \(1\pmod8\) exactly when \(V,m\) have the same parity.
There are
\(2\cdot2^{j-2}2^{j-1}=2^{2j-2}\) such pairs, and four roots \(r\),
so this subbranch gives \(2^{2j}\) triples \((r,v,m)\bmod2^j\).

Now suppose \(r=2R\) and \(v\) is odd.  The congruence in
\eqref{koide-count-eq-two-adic-odd-branch}, solved for \(v^2\), has an
odd square on the left; its right side is \(1\pmod8\) exactly when
\(R,m\) have opposite parity.  Again there are \(2^{2j-2}\) admissible
pairs and four square roots, hence \(2^{2j}\) transformed triples.

It remains to account for the divided-coordinate map.  Each class
\((r,v)\bmod2^j\) of opposite parity has exactly two inverse classes
\((a,q)\bmod2^k\): shifting either \(r\) or \(v\) by \(2^j\) changes both
\(a=r+3v\) and \(q=r+v\) by \(2^j\), while shifting both changes neither
modulo \(2^k\).  The class \(m\bmod2^j\) gives one even class
\(n=2m\bmod2^k\).  Thus the second branch contributes
\begin{equation}
 2\bigl(2^{2j}+2^{2j}\bigr)=2^{2k}.
 \label{koide-count-eq-two-adic-second-branch-count}
\end{equation}
Combining \eqref{koide-count-eq-two-adic-first-branch-count} and
\eqref{koide-count-eq-two-adic-second-branch-count} gives
\begin{equation}
 N_{2,k}=3\cdot2^{2k-1}\qquad(k\geq4),
 \label{koide-count-eq-two-adic-total-count}
\end{equation}
and hence \(\sigma_2=3/2\).

For \(p\geq5\), one has \(p\nmid108\).  Borovoi's good-prime formula
\cite[Section~1.3]{Borovoi_2001} gives
\[
 \sigma_p=\frac{1-p^{-2}}{1-\chi(p)p^{-1}},
 \qquad
 \chi(p)=\left(\frac{-(-2)\det A}{p}\right)
 =\left(\frac{108}{p}\right)=\left(\frac3p\right),
\]
which proves the remaining formula in
\eqref{koide-count-eq-all-local-densities}.  Since
\(\chi(p)=\chi_{12}(p)\) for \(p\geq5\),
\[
 \prod_{p\geq5}\sigma_p
 =L(1,\chi_{12})\prod_{p\geq5}(1-p^{-2})
 =\frac32\frac{L(1,\chi_{12})}{\zeta(2)}.
\]
Multiplication by \(\sigma_2\sigma_3=3\) gives the first equality in
\eqref{koide-count-eq-singular-series}.  Finally, the residue classes on
which \(\chi_{12}\) is nonzero give
\begin{align*}
 L(1,\chi_{12})
 &=\sum_{m\geq0}
 \left(
 \frac1{12m+1}-\frac1{12m+5}
 -\frac1{12m+7}+\frac1{12m+11}
 \right)\\
 &=\int_0^1\frac{1-x^2}{x^4-x^2+1}\,\mathrm dx
 =\frac{\log(2+\sqrt3)}{\sqrt3},
\end{align*}
where the middle equality follows by summing the geometric series and the
last by an elementary partial-fraction integration.  This completes the
local computation.
\end{proof}

\subsection{The archimedean chamber factor}
\label{koide-count-subsec-archimedean}

The Poincar\'e-residue gauge form on the whole smooth affine quadric is the
unique regular \(2\)-form \(\omega_\infty\) whose local representatives
satisfy
\begin{equation}
 \omega_\infty\wedge\mathrm dF
 =\mathrm da\wedge\mathrm dq\wedge\mathrm dn,
 \label{koide-count-eq-gauge-form}
\end{equation}
Explicitly, on the three standard charts it is
\begin{equation}
 \omega_\infty
 =\frac{\mathrm dq\wedge\mathrm dn}{2a}\quad(a\neq0),
 \qquad
 \omega_\infty
 =\frac{\mathrm da\wedge\mathrm dn}{6q}\quad(q\neq0),
 \qquad
 \omega_\infty
 =-\frac{\mathrm da\wedge\mathrm dq}{36n}\quad(n\neq0).
 \label{koide-count-eq-global-gauge-charts}
\end{equation}
These expressions agree on overlaps after restriction to the tangent bundle
of \(F=-2\).  Since the gradient
\((2a,-6q,-36n)\) cannot vanish on that level, the three charts cover
\(\mathscr X\), so \(\omega_\infty\) is globally defined and nowhere
vanishing.  Let \(m_\infty:=|\omega_\infty|\) be its associated positive
measure.  On the sheet \(a>0\), used for the chamber count, the first
expression in \eqref{koide-count-eq-global-gauge-charts} applies.
Let
\begin{equation}
 \mathcal C:=
 \{(a,q,n)\in\mathbb R^3:a>0,\ q\geq0,\ n\geq0,\ a+4n>q\}.
 \label{koide-count-eq-real-chamber}
\end{equation}

\begin{theorem}[Exact archimedean factor]
\label{koide-count-thm-archimedean-factor}
As \(X\to\infty\),
\begin{equation}
 m_\infty\{u\in\mathscr X(\mathbb R)\cap\mathcal C:
              4a+17n\leq X\}
 \sim I_\infty X,
 \qquad
 I_\infty=\frac1{2\sqrt3}
 \log\!\left(\frac3{2\sqrt2}\right).
 \label{koide-count-eq-archimedean-asymptotic}
\end{equation}
\end{theorem}

\begin{proof}
First, the wall \(a+4n>q\) removes only a bounded part of the positive
orthant.  Indeed, if \(q\geq a+4n\), then \(a\leq q-4n\); using
\(a^2=3q^2+18n^2-2\) gives
\[
 3q^2+18n^2-2\leq(q-4n)^2,
 \qquad\text{hence}\qquad
 q^2+4qn+n^2\leq1.
\]
Thus the removed real region has bounded \(m_\infty\)-measure.  Its only
integral point on \(\mathscr X\) with \(a>0,q,n\geq0\) is
\((a,q,n)=(1,1,0)\).

Introduce polar coordinates in the negative plane by
\begin{equation}
 \sqrt3\,q=r\cos\theta,
 \qquad
 3\sqrt2\,n=r\sin\theta,
 \qquad
 0\leq\theta\leq\frac\pi2.
 \label{koide-count-eq-polar-coordinates}
\end{equation}
On \(F=-2\), one has \(a=\sqrt{r^2-2}\) and \(r\,\mathrm dr=a\,\mathrm da\).
The Jacobian in \eqref{koide-count-eq-polar-coordinates} therefore gives
\begin{equation}
 \mathrm dm_\infty
 =\frac1{6\sqrt6}\,\mathrm da\,\mathrm d\theta.
 \label{koide-count-eq-polar-measure}
\end{equation}
At fixed \(\theta\), the height boundary is
\begin{equation}
 4a+\frac{17}{3\sqrt2}\sqrt{a^2+2}\sin\theta=X.
 \label{koide-count-eq-height-boundary}
\end{equation}
Let \(A_X(\theta)\) be its nonnegative solution when one exists, and put
it equal to zero otherwise.  Then \(0\leq A_X(\theta)/X\leq1/4\), and
for every fixed \(\theta\),
\[
 \frac{A_X(\theta)}{X}
 \longrightarrow
 \frac1{4+\frac{17}{3\sqrt2}\sin\theta}.
\]
Dominated convergence and \eqref{koide-count-eq-polar-measure} give
\begin{equation}
 I_\infty
 =\frac1{6\sqrt6}\int_0^{\pi/2}
 \frac{\mathrm d\theta}
 {4+\frac{17}{3\sqrt2}\sin\theta}.
 \label{koide-count-eq-archimedean-integral}
\end{equation}
For completeness, set \(t=\tan(\theta/2)\).  If
\(\alpha=2\sqrt2/3\) and \(\beta=3/(2\sqrt2)\), then
\(\alpha\beta=1\), and
\[
 \int_0^{\pi/2}
 \frac{\mathrm d\theta}
 {4+\frac{17}{3\sqrt2}\sin\theta}
 =\frac12\int_0^1\frac{\mathrm dt}{(t+\alpha)(t+\beta)}
 =3\sqrt2\log\!\left(\frac3{2\sqrt2}\right).
\]
Substitution in \eqref{koide-count-eq-archimedean-integral} proves
\eqref{koide-count-eq-archimedean-asymptotic}.
\end{proof}

\subsection{A genuine norm and the homogeneous counting theorem}
\label{koide-count-subsec-gos}

\begin{lemma}[A genuine norm realizing the height]
\label{koide-count-lem-genuine-norm}
Define
\begin{equation}
 \lVert(a,q,n)\rVert_*
 :=\max\{|4a+17n|,|q|,|a|\}.
 \label{koide-count-eq-star-norm}
\end{equation}
Then \(\lVert\cdot\rVert_*\) is a norm on \(\mathbb R^3\), and on the
chamber \(\mathcal C\) it satisfies
\begin{equation}
 \lVert(a,q,n)\rVert_*=4a+17n.
 \label{koide-count-eq-norm-is-height}
\end{equation}
For every integer \(X\geq0\),
\begin{equation}
 R(X)=\#\{u\in\mathscr X(\mathbb Z)\cap\mathcal C:
                    \lVert u\rVert_*<X+1\}.
 \label{koide-count-eq-r-as-sector-count}
\end{equation}
\end{lemma}

\begin{proof}
The linear map
\[
 (a,q,n)\longmapsto(4a+17n,q,a)
\]
has determinant \(-17\).  Thus \eqref{koide-count-eq-star-norm} is the
pullback of the supremum norm by an invertible linear map and is a genuine
norm.  If \((a,q,n)\in\mathcal C\), then
\[
 0<a<4a+17n,
 \qquad
 0\leq q<a+4n<4a+17n,
\]
which proves \eqref{koide-count-eq-norm-is-height}.  Since \(4a+17n\)
is integral on \(\mathscr X(\mathbb Z)\), equation
\eqref{koide-count-eq-r-as-sector-count} follows from
\eqref{koide-count-eq-exact-count}.
\end{proof}

\begin{lemma}[Verification of the Gorodnik--Oh--Shah hypotheses]
\label{koide-count-lem-gos-hypotheses}
Let \(\mathbf G=\operatorname{Spin}(F)\), acting through its standard
three-dimensional orthogonal representation on \(W=\mathbb R^3\), and let
\(x_0=(1,1,0)\).  Then
\[
 \mathscr X(\mathbb R)=\mathbf G(\mathbb R)^\circ x_0.
\]
This orbit is a closed affine symmetric variety; \(\mathbf G\) is connected
and absolutely almost simple over \(\mathbb Q\), is isotropic over
\(\mathbb R\), and acts irreducibly on \(W\) over \(\mathbb R\); and the
stabilizer \(\mathbf H=\operatorname{Stab}_{\mathbf G}(x_0)\) has no
nontrivial \(\mathbb Q\)-characters.  The cone \(\mathcal C\) in
\eqref{koide-count-eq-real-chamber} is admissible and generic for
\cite[Theorem~1.6]{GorodnikOhShah_2009}.
\end{lemma}

\begin{proof}
The form \(F\) has signature \((1,2)\).  Hence \(\mathbf G\) is a
connected absolutely almost simple group of type \(A_1\), it is
\(\mathbb R\)-isotropic, and its standard three-dimensional representation
is \(\mathbb R\)-irreducible.  The nonzero level \(\mathscr X\) is
Zariski closed.  With
\[
 Q=\sqrt3\,q,
 \qquad
 N=3\sqrt2\,n,
\]
its real equation becomes
\[
 Q^2+N^2-a^2=2.
\]
The parametrization
\[
 (a,\theta)\longmapsto
 \bigl(a,\sqrt{a^2+2}\cos\theta,
          \sqrt{a^2+2}\sin\theta\bigr),
 \qquad (a,\theta)\in\mathbb R\times(\mathbb R/2\pi\mathbb Z),
\]
shows that \(\mathscr X(\mathbb R)\) is connected.  Witt extension first
produces an element \(t\in\operatorname O(F)(\mathbb R)\) taking any
prescribed point of the level set to any other.  If \(\det t=-1\), choose
an anisotropic vector in the orthogonal complement of the target point and
compose \(t\) on the left with its reflection.  That reflection fixes the
target point and has determinant \(-1\); hence the composite lies in
\(\operatorname{SO}(F)(\mathbb R)\).  Thus
\(\operatorname{SO}(F)(\mathbb R)\) is transitive on the level set.  The
identity-component orbits have the full dimension of the level surface and
are open; because their complements are unions of open orbits, each is also
closed.  Connectedness therefore leaves only one such orbit, proving
\(\mathscr X(\mathbb R)=\mathbf G(\mathbb R)^\circ x_0\).  The stabilizer
of a nonisotropic vector is the fixed subgroup of the involution induced by
reflection in that vector, so the orbit is affine symmetric.

The orthogonal complement of \(x_0\) is
\begin{equation}
 x_0^\perp=\{(3s,s,t):s,t\in\mathbb Q\},
 \qquad
 F(3s,s,t)=6(s^2-3t^2).
 \label{koide-count-eq-rational-complement}
\end{equation}
Thus \(\mathbf H\) is the norm-one torus of
\(K=\mathbb Q(\sqrt3)\).  Its character lattice is \(\mathbb Z\), with
the nontrivial element of \(\operatorname{Gal}(K/\mathbb Q)\) acting by
\(-1\); consequently it has no nonzero Galois-invariant character and no
nontrivial \(\mathbb Q\)-character.  The integral point \(x_0\) supplies
the required nonempty integral locus.

It remains to check the sector rather than merely the affine orbit.  For the
given representation and any Euclidean sphere in \(W\), the Satake boundary
is
\begin{equation}
 \mathscr X^\infty=\{F=0\}\cap S(W)
 =\mathcal O_+\sqcup\mathcal O_-,
 \qquad
 \mathcal O_\pm:=\{F=0,\ \pm a>0\}\cap S(W),
 \label{koide-count-eq-satake-boundary}
\end{equation}
and \(\mathbf G(\mathbb R)^\circ\) is transitive on each displayed circle.
For a nonzero null vector with \(a,q,n\geq0\),
\[
 a=\sqrt{3q^2+18n^2}>q-4n.
\]
Hence both the closure and the relative interior of \(\mathcal C\) meet
exactly \(\mathcal O_+\), not \(\mathcal O_-\).  The intersection of
\(\partial\mathcal C\) with the Satake boundary consists of the two radial
directions
\begin{equation}
 (a,q,n)=(\sqrt3,1,0),
 \qquad
 (a,q,n)=(3\sqrt2,0,1).
 \label{koide-count-eq-sector-boundary-rays}
\end{equation}
These are exactly the two projective limits of the integral Pell families
in Proposition~\ref{koide-count-prop-pell-faces}.
This two-point intersection has measure zero for the smooth measure class
on the boundary orbit.
This is precisely the admissibility and genericity condition required by
the cited sector theorem.
\end{proof}

\subsection{The 2-adic orbit sign and Borovoi's coefficient}
\label{koide-count-subsec-borovoi}

Write $\widehat{\mathbb Z}=\prod_p\mathbb Z_p$.  For the nonsquare
ternary quadric under consideration, Borovoi associates to every local
point an orbit sign $\nu_v\in\{\pm1\}$, constant on
$\operatorname{Spin}(F)(\mathbb Q_v)$-orbits and equal to $+1$ at almost
every finite place.  Thus
$\nu_f((x_p)_p)=\prod_p\nu_p(x_p)$ is defined on
$\mathscr X(\widehat{\mathbb Z})$.  Because $\mathscr X(\mathbb R)$ is
connected, his normalization reads
\begin{equation}
 c_{\mathscr X}
 =\frac{2m_f\{x\in\mathscr X(\widehat{\mathbb Z}):\nu_f(x)=+1\}}
 {m_f(\mathscr X(\widehat{\mathbb Z}))},
 \label{koide-count-eq-borovoi-sign-definition}
\end{equation}
where $m_f$ is the finite Tamagawa measure attached to the gauge form
\eqref{koide-count-eq-gauge-form} and Borovoi's convergence factors.

\begin{proposition}[Finite adelic sign balance]
\label{koide-count-prop-finite-sign-balance}
For the affine quadric \(\mathscr X=\{F=-2\}\), the coefficient
\(c_{\mathscr X}\) in Borovoi's counting formula
\cite[Theorem~1.4.1 and Proposition~1.5.1]{Borovoi_2001} is exactly
\begin{equation}
 c_{\mathscr X}=1.
 \label{koide-count-eq-borovoi-coefficient}
\end{equation}
\end{proposition}

\begin{proof}
Continue to write \(A=\operatorname{diag}(1,-3,-18)\).  Put
\(e=(1,0,0)\), \(v=(1,0,1)\), and for every anisotropic vector \(w\) let
\[
 r_w(u)=u-\frac{2u^{\mathsf T}Aw}{w^{\mathsf T}Aw}\,w
\]
be its orthogonal reflection.  The element \(g=r_vr_e\) is
\begin{equation}
 g=
 \begin{pmatrix}
  -19/17&0&-36/17\\
  0&1&0\\
  -2/17&0&-19/17
 \end{pmatrix}.
 \label{koide-count-eq-sign-toggle-matrix}
\end{equation}
Direct multiplication gives
\begin{equation}
 g^{\mathsf T}Ag=A,
 \qquad
 \det g=1.
 \label{koide-count-eq-sign-toggle-orthogonal}
\end{equation}
Because \(17\in\mathbb Z_2^\times\), both \(g\) and \(g^{-1}\) have entries
in \(\mathbb Z_2\), so \(g\in\operatorname{SO}(F,\mathbb Z_2)\).  Its
spinor norm is
\begin{equation}
 \vartheta(g)=F(v)F(e)=-17
 \quad\text{in }\mathbb Q_2^\times/\mathbb Q_2^{\times2}.
 \label{koide-count-eq-toggle-spinor-norm}
\end{equation}

We identify the stabilizer's spinor-norm image.  By
\eqref{koide-count-eq-rational-complement}, over \(\mathbb Q_2\) the
orthogonal complement of \(x_0\) is
\begin{equation}
 F|_{x_0^\perp}
 =6(s^2-3t^2)
 =6N_{K_2/\mathbb Q_2}(s+t\sqrt3),
 \qquad K_2=\mathbb Q_2(\sqrt3).
 \label{koide-count-eq-two-adic-complement}
\end{equation}
The odd unit \(3\) is not a square in \(\mathbb Q_2\), so \(K_2\) is a
quadratic field and this binary plane is anisotropic.  Put
\(u_0=(3,1,0)\), so \(F(u_0)=6\).  If
\(h\in\operatorname{SO}(F|_{x_0^\perp})(\mathbb Q_2)\), then
\(hr_{u_0}\) is a two-dimensional orthogonal transformation of determinant
\(-1\), hence is a reflection \(r_w\) in some nonisotropic vector
\(w=(3s,s,t)\).  Therefore \(h=r_wr_{u_0}\) and
\[
 \vartheta(h)=F(w)F(u_0)
 =36N_{K_2/\mathbb Q_2}(s+t\sqrt3)
 \quad\bmod\mathbb Q_2^{\times2}.
\]
Conversely these products realize every norm square class as \(w\) ranges
through the nonzero vectors of the anisotropic plane.  Hence
\begin{equation}
 \vartheta\!\left(
 \operatorname{Stab}_{\operatorname{SO}(F)}(x_0)(\mathbb Q_2)
 \right)
 =N_{K_2/\mathbb Q_2}(K_2^\times)
 \quad\bmod\mathbb Q_2^{\times2}.
 \label{koide-count-eq-stabilizer-spinor-image}
\end{equation}

For a quadratic extension, its norm group is the kernel of the Hilbert
symbol with the defining square class
\cite[Corollary~13.4.6]{Voight_2021}.  The odd-unit formula
\cite[Equation~12.4.14]{Voight_2021} gives
\begin{equation}
 (-17,3)_2
 =(-1)^{((-17-1)/2)((3-1)/2)}=-1.
 \label{koide-count-eq-minus17-hilbert-symbol}
\end{equation}
Thus \(-17\notin N_{K_2/\mathbb Q_2}(K_2^\times)\).

By Witt extension, every \(x\in\mathscr X(\mathbb Q_2)\) is \(tx_0\) for
some \(t\in\operatorname O(F)(\mathbb Q_2)\).  If \(\det t=-1\), choose
an anisotropic vector in \(x^\perp\) and compose on the left with its
reflection; this fixes \(x\) and changes the determinant.  Hence every
such \(x\) has the form \(x=sx_0\) with
\(s\in\operatorname{SO}(F)(\mathbb Q_2)\).  The image of
\(\operatorname{Spin}(F)(\mathbb Q_2)\) in
\(\operatorname{SO}(F)(\mathbb Q_2)\) is the kernel of the spinor norm,
as used in \cite[Lemma~3.1]{Borovoi_2001}.  Consequently the
\(\operatorname{Spin}(F)(\mathbb Q_2)\)-orbit of \(x=sx_0\) is determined
by
\begin{equation}
 \vartheta(s)\in
 \mathbb Q_2^\times/
 \bigl(\mathbb Q_2^{\times2}
 N_{K_2/\mathbb Q_2}(K_2^\times)\bigr),
 \label{koide-count-eq-two-adic-orbit-quotient}
\end{equation}
a quotient of order two.  Equations
\eqref{koide-count-eq-toggle-spinor-norm} and
\eqref{koide-count-eq-minus17-hilbert-symbol} show that multiplication by
\(g\) exchanges these two orbits.  In Borovoi's convention, the orbit of
\(x_0\) has sign \(\nu_2=+1\) and the other orbit has sign
\(\nu_2=-1\).  Hence
\begin{equation}
 \nu_2(gx)=-\nu_2(x)
 \qquad(x\in\mathscr X(\mathbb Q_2)).
 \label{koide-count-eq-local-sign-toggle}
\end{equation}

Apply \(g\) at the prime \(2\) and the identity at every other finite
prime.  This gives a bijection of \(\mathscr X(\widehat{\mathbb Z})\).  It
preserves the finite Tamagawa measure: at \(2\) this follows from invariance
of the local measure, equivalently from \(g^*F=F\) and \(\det g=1\), and at
all other primes the map is the identity.  Equation
\eqref{koide-count-eq-local-sign-toggle} shows that it exchanges the finite
sign loci \(\nu_f=+1\) and \(\nu_f=-1\).  They therefore have equal
measure.  Since \(\mathscr X(\mathbb R)\) is connected, Borovoi's formula
becomes
\[
 c_{\mathscr X}
 =\frac{2m_f(\mathscr X(\widehat{\mathbb Z})_+)}
        {m_f(\mathscr X(\widehat{\mathbb Z}))}=1.
\]
The argument uses a measure-preserving, sign-reversing bijection; it does
not require \(g\) to be an involution (indeed \(g^2\neq I\)).
\end{proof}

\subsection{Calibration and the counting constant}
\label{koide-count-subsec-calibration}

\begin{theorem}[Calibrated reduced-root asymptotic]
\label{koide-count-thm-linear-asymptotic}
The number of reduced positive defect-one roots satisfies
\begin{equation}
 R(X)\sim C_0X\qquad(X\to\infty),
 \label{koide-count-eq-final-asymptotic}
\end{equation}
where
\begin{equation}
 \begin{aligned}
 C_0
 &=\mathfrak S(\mathscr X)I_\infty\\
 &=\frac9{2\pi^2}\log(2+\sqrt3)
   \log\!\left(\frac3{2\sqrt2}\right)\\
 &=0.035362047799731853\ldots.
 \end{aligned}
 \label{koide-count-eq-final-constant}
\end{equation}
\end{theorem}

\begin{proof}
Let \(\mu_{\mathrm{GOS}}\) denote the invariant measure on
\(\mathscr X(\mathbb R)\) with the normalization appearing in the
Gorodnik--Oh--Shah sector theorem
\cite[Theorem~1.6]{GorodnikOhShah_2009}.  The gauge measure \(m_\infty\)
in \eqref{koide-count-eq-gauge-form} is also a nonzero
\(\mathbf G(\mathbb R)^\circ\)-invariant Radon measure.  Because
\(\mathscr X(\mathbb R)\) is a single homogeneous orbit, uniqueness of
invariant measure gives
\begin{equation}
 \mu_{\mathrm{GOS}}=\kappa m_\infty
 \label{koide-count-eq-measure-scalar}
\end{equation}
for some \(\kappa>0\).

The scalar is fixed by applying both counting theorems to the same Euclidean
balls and the full cone \(W\).  Gorodnik--Oh--Shah give
\[
 \#\{x\in\mathscr X(\mathbb Z):|x|_2<T\}
 \sim\kappa\,m_\infty
 \{x\in\mathscr X(\mathbb R):|x|_2<T\}.
\]
Borovoi's Theorem~1.4.1 and Proposition~1.5.1, together with
Proposition~\ref{koide-count-prop-finite-sign-balance}, give for the same
balls
\[
 \#\{x\in\mathscr X(\mathbb Z):|x|_2<T\}
 \sim\mathfrak S(\mathscr X)\,m_\infty
 \{x\in\mathscr X(\mathbb R):|x|_2<T\}.
\]
The real volume here is asymptotic to a positive constant times \(T\).
Comparison therefore gives the invariant-measure calibration
\begin{equation}
 \kappa=\mathfrak S(\mathscr X).
 \label{koide-count-eq-measure-calibration}
\end{equation}

Lemma~\ref{koide-count-lem-genuine-norm} and
Lemma~\ref{koide-count-lem-gos-hypotheses} now allow the sector theorem to
be applied with the actual norm and chamber of the reduced-root problem:
\begin{align*}
 R(X)
 &\sim\mu_{\mathrm{GOS}}
 \{u\in\mathscr X(\mathbb R)\cap\mathcal C:
       \lVert u\rVert_*\leq X\}\\
 &=\mathfrak S(\mathscr X)\,
 m_\infty\{u\in\mathscr X(\mathbb R)\cap\mathcal C:
       4a+17n\leq X\}\\
 &\sim\mathfrak S(\mathscr X)I_\infty X.
\end{align*}
The switch between strict and weak inequalities is harmless by a
\((1\pm\varepsilon)X\) squeeze; on integral points it also follows directly
from \eqref{koide-count-eq-r-as-sector-count}.  Since
\(R(X)=R(\lfloor X\rfloor)\), the integer-height asymptotic extends to all
real \(X\to\infty\).  Finally, substituting
\eqref{koide-count-eq-singular-series} and
\eqref{koide-count-eq-archimedean-asymptotic}, and using
\(\zeta(2)=\pi^2/6\), gives exactly
\eqref{koide-count-eq-final-constant}.
\end{proof}

\subsection{Exact computation and the remaining open problem}
\label{koide-count-subsec-computation}

The first five reduced roots are
\begin{equation}
 (1,4,20),\quad(4,4,33),\quad(4,8,49),\quad
 (4,15,76),\quad(8,24,129).
 \label{koide-count-eq-first-roots}
\end{equation}
The first root is the first admissible term of the $n=0$ Pell family and
the second is the first term of the $q=0$ Pell family in
Proposition~\ref{koide-count-prop-pell-faces}.
Exact integer enumeration in the ternary coordinates gives the following
values.  The residual column uses the full constant in
\eqref{koide-count-eq-final-constant}, not a fitted constant.
\begin{center}
\begin{tabular}{r@{\qquad}r@{\qquad}r}
\toprule
\(X\)&\(R(X)\)&\(R(X)-C_0X\)\\
\midrule
10,000&353&\(-0.62048\)\\
20,000&715&\(7.75904\)\\
50,000&1,755&\(-13.10239\)\\
75,000&2,654&\(1.84642\)\\
100,000&3,535&\(-1.20478\)\\
125,000&4,460&\(39.74403\)\\
150,000&5,294&\(-10.30717\)\\
175,000&6,195&\(6.64164\)\\
200,000&7,109&\(36.59044\)\\
\bottomrule
\end{tabular}
\end{center}
The enumeration loops over \(n\geq0\) and \(a>0\) with
\(4a+17n\leq X\), tests whether
\[
 q^2=\frac{a^2-18n^2+2}{3}
\]
is the square of a nonnegative integer satisfying \(q<a+4n\), and then
reconstructs \((x,y,z)\) with
\eqref{koide-count-eq-inverse-coordinates}.  Every retained triple is
rechecked in the original polynomial \eqref{koide-count-eq-defect} and
against all ordering and wall inequalities.  Separate exhaustive congruence
enumerations agree with the formulae above at
\(p=2,3,5,7,11,13\) through the computed powers.  These finite checks
verify the implementation; the asymptotic theorem rests on the proofs above,
not on numerical fitting.

The reproducibility script is
\begin{center}
 \path{supporting_materials/workbench/reproduction/t1_koide_exact/verify_koide_narain_root_density.py}.
\end{center}
Running it with \texttt{--max-z 200000} emits the coordinate, local-density,
$2$-adic isometry, and bounded-enumeration certificates recorded in
\begin{center}
 \path{supporting_materials/workbench/reproduction/t1_koide_exact/koide_verification_200000.json}.
\end{center}
The script uses vectorized square detection only to propose candidates; it
rechecks every retained candidate with integer arithmetic in both coordinate
systems and in the original polynomial $\Delta$.

\begin{test}[Square-root discrepancy]
\label{koide-count-open-square-root-discrepancy}
\Open Determine whether, for every \(\varepsilon>0\),
\begin{equation}
 R(X)-C_0X=O_\varepsilon(X^{1/2+\varepsilon}).
 \label{koide-count-eq-square-root-discrepancy}
\end{equation}
The preceding homogeneous-dynamics argument proves
\(R(X)=C_0X+o(X)\), but does not prove
\eqref{koide-count-eq-square-root-discrepancy}.  No unproved spectral gap,
GRH assumption, or implicit ``conditional theorem'' is being used here:
the sharper error term is an open problem in this workbench.
\end{test}
